\subsection{\texorpdfstring{关系式 \(\sum_{i} e_{i} f_{i} = n\)}{关系式 Σ ei fi = n}}

命题 5. 设 K 为域，v 为 K 上的赋值，A 为其环，m 为其理想，L 为 K 的次数为 n 的有限扩张，B 为 A 在 L 中的整闭包，且 \((v_{i}^{\prime})_{1 \leq i \leq s}\) 是 v 到 L 的完备延拓系。则

\[
[ \mathrm{B} / \mathrm{mB}: \mathrm{A} / \mathrm{m} ] = \sum_ {i = 1} ^ {s} \varepsilon (v _ {i} ^ {\prime} / v) f (v _ {i} ^ {\prime} / v).
\]

令 \(A_{i}\) 为 \(v_{i}^{\prime}\) 的环；则 \(A_{i} = B_{m_{i}}\)，其中 \(m_{i}\) 遍历 B 的极大理想族（第 3 节，注）。令 \(q_{i}\) 为关于 \(m_{i}\) 的 mB 的饱和（第二章，第 2 节，第 4 节）。根据第五章第 2 节第 1 节命题 1 的推论 3，典范同态 \(B/mB \rightarrow \prod_{i=1}^{s} B/q_i\) 是同构，并且 \(m_i\) 是 B 中包含 \(q_i\) 的唯一极大理想。因此，\(B/q_i\) 典范同构于 \((\mathrm{B}/q_i)_{m_i}\)（第二章，第 3 节，第 3 节，命题 8），也就是

\[
\mathrm{B} _ {\mathrm{m} _ {i}} / \mathrm{mB} _ {\mathrm{m} _ {i}} = \mathrm{A} _ {i} / \mathrm{mA} _ {i}.
\]

因此存在典范同构 \(B/mB \rightarrow \prod_{i=1}^{s} A_{i}/mA_{i}\)，由第 4 节命题 4 即得结果。

推论。在相同假设和记号下，

\[
[ \mathrm{B} / \mathrm{mB}: \mathrm{A} / \mathrm{m} ] = \sum_ {i = 1} ^ {s} \varepsilon (v _ {i} ^ {\prime} / v) f (v _ {i} ^ {\prime} / v) \leqslant \sum_ {i = 1} ^ {s} e (v _ {i} ^ {\prime} / v) f (v _ {i} ^ {\prime} / v) \leqslant n.
\]

我们知道 \(\varepsilon(v_i'/v) \leqslant e(v_i'/v)\)（第 4 节，命题 3 的推论），并且 \(\sum_{i=1}^{s} e(v_i'/v)f(v_i'/v) \leqslant n\)（第 3 节，定理 1）。

定理 2. 在命题 5 的假设和记号下，下列条件等价：

\begin{enumerate}
\def\labelenumi{(\alph{enumi})}
\item
  B 是有限生成的 A-模；
\item
  B 是自由 A-模；
\item
  \([\mathbf{B} / \mathfrak{m}\mathbf{B}\colon \mathbf{A} / \mathfrak{m}] = n;\)
\end{enumerate}

\[
\text {(d)} \sum_ {i = 1} ^ {n} e (v _ {i} ^ {\prime} / v) f (v _ {i} ^ {\prime} / v) = n \text { and } \varepsilon (v _ {i} ^ {\prime} / v) = e (v _ {i} ^ {\prime} / v) \text { for   all } i.
\]

条件 (a) 与 (b) 的等价性由第 3 节第 6 节引理 1 得到。显然 (b) 蕴含 (c)（代数，第二章，第 1 节，第 5 节，公式 (19)）。(c) 与 (d) 的等价性由命题 5 的推论得到。还需证明 (c) 蕴含 (b)。

一般地，若 M 是 A-模，我们把 M/mM 关于 A/m 的向量空间记为 V(M)。假设 (c) 意味着 \(\dim(\mathrm{V}(\mathrm{B})) = n\)。令 \(x_{1}, \ldots, x_{n}\) 为 B 中的元素，它们在 V(B) 中的典范像构成 V(B) 的一个基，并令 \(L \subset B\) 为它们生成的 A-子模。由于 L 是无挠且有限生成的，它是自由的（§ 3，第 6 节，引理 1）。我们将看到 B = L。令 \(y \in B\)；记 \(M = L + Ay\)；这也是自由 A-模。典范嵌入 \(L \to M \to B\) 给出典范同态 \(\mathrm{V}(\mathrm{L}) \to \mathrm{V}(\mathrm{M}) \to \mathrm{V}(\mathrm{B})\)。由于 L 和 M 的秩均 \(\leqslant n\)，V(L) 和 V(M) 的维数也均不超过 n。现在由假设，\(\mathrm{V}(\mathrm{L}) \to \mathrm{V}(\mathrm{B})\) 是满射，而 V(B) 是 n 维的，故 V(L) 和 V(M) 都是 n 维的，且 \(\mathrm{V}(\mathrm{L}) \to \mathrm{V}(\mathrm{M})\) 是满射。由于 M 是有限生成的，\(L \to M\) 是满射（第二章，第 3 节，第 2 节，命题 4 的推论 1），从而 \(L = M, y \in L\)，且 B = L。因此 B 是自由的。

注 (1)。若 \(v\) 离散，\(\varepsilon(v_i'/v) = e(v_i'/v)\)（第 4 节），条件 (d) 化为 \(\sum_{i=1}^{s} e(v_i'/v)f(v_i'/v) = n\)。

推论 1. 在相同假设和记号下，再设 v 离散且 L 可分。则

\[
\sum_ {i = 1} ^ {s} e (v _ {i} ^ {\prime} / v) f (v _ {i} ^ {\prime} / v) = n.
\]

此时 A 的整闭包 B 是秩 n 的自由 A-模，因为 A 是主理想整环（第五章，第 1 节，第 6 节，命题 18 的推论 2）。

推论 2. 设 K 为域，v 是 K 上使 K 完备的离散赋值，L 为 K 的次数为 n 的有限扩张。~则 v 到 L 存在唯一（在等价意义下）的延拓 v'，v' 的环 A' 是 v 的环 A 上的有限生成自由模，且 \(e(v'/v)f(v'/v)=n\)。

\(v\) 到 \(L\) 的所有延拓都是相关的（第 2 节，命题 2 的推论）；由于它们是离散的（第 1 节，命题 1 的推论 3），所以它们等价（§ 4，第 5 节，命题 6 (c)）。这说明 \(v'\) 的唯一性。\(A\) 在 \(L\) 中的整闭包于是为 \(A'\)（§ 1，第 3 节，定理 3 的推论 3）。由于 \(v\) 离散，拓扑在 \(A\) 上由 \(K\) 上的拓扑诱导，是 m-进拓扑（其中 \(m = m(A)\)）；环 \(A\) 是完备的，因为它在 \(K\) 中闭。由于 \(A'/mA'\) 是有限维的向量 \((A/m)\)-空间（第 4 节，命题 4），\(A'\) 是有限生成的 A-模（第三章，第 2 节，第 9 节，命题 12 的推论 3）。因此由定理 2，A' 是自由的且 \(e(v'/v)f(v'/v) = n\)。

推论 3. 假设 \(v\) 的高为 1，并且定理 2 的等价条件成立；若 \(\hat{\mathbf{L}}_i\) 是 \(\mathbf{L}\) 关于 \(v_i'\) 的完备化，则次数 \(n_i = [\hat{\mathbf{L}}_i: \hat{\mathbf{K}}]\) 等于 \(e(v_i'|v)f(v_i'|v)\)，对所有 \(i\) 都成立，且典范同态

\[
\phi \colon \hat {\mathbf {K}} \otimes_ {\mathbf {K}} \mathbf {L} \rightarrow \prod_ {i = 1} ^ {s} \hat {\mathbf {L}} _ {i}
\]

（第 2 节，命题 2）。是双射的。对于所有 \(x \in \mathbf{L}\)，特征多项式 \(\mathrm{Pc}_{\mathrm{L} / \mathbb{K}}(x; \mathbf{X})\) 等于特征多项式 \(\mathrm{Pc}_{\hat{\mathrm{L}}_i / \hat{\mathbf{K}}}(x; \mathbf{X})\) 的乘积（\(1 \leqslant i \leqslant s\)）；特别地，

\[
\left\{ \begin{array} { l } \mathrm{Tr} _ { \mathbf { L } / \mathbb { K } } ( x ) = \sum _ { i = 1 } ^ { s } \mathrm{Tr} _ { \hat { \mathbf { L } } _ { i } / \hat { \mathbf { K } } } ( x ) \\ \mathrm{N} _ { \mathbf { L } / \mathbb { K } } ( x ) = \prod _ { i = 1 } ^ { s } \mathrm{N} _ { \hat { \mathbf { L } } _ { i } / \hat { \mathbf { K } } } ( x ) \\ v ( \mathrm{N} _ { \mathbf { L } / \mathbb { K } } ( x ) ) = \sum _ { i = 1 } ^ { s } n _ { i } v _ { i } ^ { \prime } ( x ) . \\ \end{array} \right.\tag{9}
\]

（(9) 中最后一个关系是有意义的，因为显然可以假设 \(v_{i}^{\prime}\)（由第 1 节命题 1 的推论 2 知其高为 1）像 \(v\) 一样将其值取在 \(\mathbf{R}\) 的一个子群中。）

由于 \(v_{i}^{\prime}\) 两两不等价且它们的高为 1，它们独立，因此第 2 节命题 2 表明 \(e(v_{i}^{\prime}/v)f(v_{i}^{\prime}/v) \leqslant n_{i}\) 对所有 \(i\) 成立，且 \(\sum_{i=1}^{s} n_i \leqslant n\)。所以第一个断言由这些不等式以及关系 \(\sum_{i=1}^{s} e(v_i'/v) f(v_i'/v) = n\) 得到。在同构 \(\phi\) 下，自同态 \(z \mapsto z(1 \otimes x)\) 作用于 \(\hat{\mathbf{K}} \otimes_{\mathbb{K}} \mathbf{L}\)（其中 \(x \in \mathbf{L}\)）时，变为 \(\prod_{i=1}^{s}\hat{L}_i\) 上的自同态；它保持每个因子不变，并在每个因子上化为乘以 \(x\)（因为 L 典范嵌入其完备化 \(\hat{L}_i\)）；由此得到关于 \(x\) 的特征多项式以及 (9) 的前两个公式的断言。最后，令 E 为 \(\hat{K}\) 的一个有限拟 Galois 扩张，包含 \(\hat{L}_i\)；由于 \(\hat{K}\) 完备且 \(\hat{v}\) 的高为 1，存在唯一一个作用于 E 的赋值 \(w\)（在等价意义下），它延拓 \(\hat{v}\)（第 2 节，命题 2 的推论 1）；于是对 E 的每个 \(\hat{K}\)-自同构 \(\sigma\)，都有 \(w(\sigma(x)) = v_i'(x)\)。因此

\[
\hat {v} \left(\mathrm{N} _ {\hat {\mathrm{L}} _ {i} / \hat {\mathrm{K}}} (x)\right) = n _ {i} v _ {i} ^ {\prime} (x)
\]

（代数，第 VIII 章，第 12 节，第 2 节，公式 (15)），这证明了 (9) 中的公式。

推论 4. 在推论 3 的假设下，若 L 是 K 的可分扩张，则每个 \(\hat{\mathbf{L}}_i\) 都是 \(\hat{\mathbf{K}}\) 的可分扩张。若进一步 L 是 K 的 Galois 扩张，Galois 群为 \(\mathcal{G}\)，且 \(\mathcal{G}_i\) 表示 B 中 \(v_i'\) 的理想的分解群（第五章，第 2 节，第 2 节，定义 2），则 \(\hat{\mathbf{L}}_i\) 是 \(\hat{\mathbf{K}}\) 的 Galois 扩张，其 Galois 群同构于 \(\mathcal{G}_i\)。

显然 \(\hat{\mathbf{L}}_i = \hat{\mathbf{K}}(\mathbf{L})\)；因此，若 \(\mathbf{L}\) 是 \(\mathbf{K}\) 的可分扩张，则 \(\hat{\mathbf{L}}_i\) 是 \(\hat{\mathbf{K}}\) 的可分扩张（代数，第五章，第 7 节，第 6 节，命题 10）。现在设 \(\mathbf{L}\) 是 Galois 扩张。每个 \(\sigma \in \mathcal{G}_i\) 都是 \(\mathbf{L}\) 上关于由 \(v_i'\) 定义的拓扑的连续自同构；由于 \(v_i'\) 的任意两个理想关于包含关系都不可比较（§ 7，第 2 节，定理 1 的推论 1），由此必有 \(v_i' = v_i' \circ \sigma\)，这由 \(\mathcal{G}_i\) 的定义得到；因此 \(\sigma\) 可通过连续性延拓为 \(\hat{\mathbf{K}}\)-自同构 \(\hat{\sigma}\)，作用于 \(\hat{\mathbf{L}}_i\)。这证明了 \(\hat{\mathbf{K}}\)-自同构在 \(\hat{\mathbf{L}}_i\) 上的数目至少等于 \(\operatorname{Card}(\mathcal{G}_i)\)。但由于赋值 \(v_i'\) 在 \(\mathcal{G}\) 下两两共轭（第五章，第 2 节，第 3 节，命题 6），有 \(s = (\mathcal{G}: \mathcal{G}_i)\)，从而

\[
\operatorname{Card} \left(\mathcal {G} _ {i}\right) = n / s \leqslant n.
\]

另一方面，根据推论 3，\(n = sn_{i}\)；这证明了 \(\hat{\mathbf{L}}_{i}\) 是 \(\hat{\mathbf{K}}\) 的 Galois 扩张，并且 \(\sigma \in \mathcal{G}_{i}\) 的自同构的连续延拓恰好构成 \(\hat{\mathbf{K}}\)-自同构在 \(\hat{\mathbf{L}}_{i}\) 上的全体。

注 (2)。上述部分结果可推广到 K 为不必交换的域上的赋值情形（参见~§ 3，第 1 节）。设 L 是 K 的扩域，且 \(v'\) 是 L 上的赋值，\(v\) 是其在 K 上的限制，A' 与 A 分别是赋值 \(v'\) 与 v 的环；则可如第 1 节定义分歧指数 \(e(v'/v)\)；另一方面，\(\kappa(\mathbf{A})\) 可视为 \(\kappa(\mathbf{A}')\) 的子域，并将 \(v'\) 关于 v 的（左）剩余秩定义为数 \(f(v'/v)\)，它等于左向量 \(\kappa(\mathbf{A})\)-空间 \(\kappa(\mathbf{A}')\) 的维数；若此维数有限，则如此定义，否则取 \(+\infty\)。于是，若 L 是关于 K 的有限维左向量空间，第 1 节引理 2 及其证明保持不变。此外，若 K 关于 v 完备，则第 5 节定理 2 的推论 2 的断言（\(v'\) 的存在性除外）也成立（n 表示 L 作为关于 K 的左向量空间的维数），证明如下：

首先，\(v'\) 在 \(L\) 上定义的拓扑是豪斯多夫的，并且与其左 \(K\)-向量空间结构相容，因此 \(v\) 到 \(L\) 的两个延拓在 \(L\) 上给出相同的拓扑（§ 5，第 2 节，命题 4），这证明这些延拓至多相差一个等价关系（§ 6，第 2 节）。下面证明：若 \(m = m(A)\)，则 \(A'/mA'\) 是左向量 \((A/m)\)-空间，其维数为 \(e(v'/v)f(v'/v)\)。记 \(e = e(v'/v)\)，可设 \(v(K^*) = Z\) 且 \(v'(L^*) = e^{-1}Z\)；令 \(u'\) 是 \(L\) 中满足 \(v'(u') = e^{-1}\) 的元素，令 \(u\) 是 \(K\) 中满足 \(v(u) = 1\) 的元素；于是 \(u = zu'^e\)，其中 \(z \in L\) 且 \(v'(z) = 0\)。由于 \(m\) 由 \(u\) 生成（作为 \(A\) 的左理想或右理想），\(mA' = u'^eA' = A'u'^e\)，只需证明：对于 \(0 \leqslant k \leqslant e - 1\)，\(A'u'^k/A'u'^{k+1}\) 是左向量 \((A/m)\)-空间，维数为 \(f(v'/v)\)。但 \(t \mapsto tu'^k\) 是从左 A-模 \(A'\) 到左 A-模 \(A'u'^k\) 的同构，将 \(A'u'\) 映到 \(A'u'^{k+1}\)，因而取商给出一个 \((A/m)\)-同构，将 \(A'/A'u'\) 映到 \(A'u'^k/A'u'^{k+1}\)；由 \(f(v'/v)\) 的定义以及 \(u'\) 生成 \(A'\) 的极大理想，便得到所述断言。证明与 \(K\)、\(L\) 交换时相同（有限生成无挠 A-模为自由模这一事实按 § 3，第 6 节，引理 1 的方式证明）。

